\scanpage{17}
\mspage{15}
\[
(3.\unsure{15}.1)\qquad I'_t=\prod_{\ell\ne p}(I'/(I',I'))(\ell)\simeq\prod_{\ell\ne p}\mathbb Z_\ell(1).
\]
\deleted{Ceci posé, on constate aussitôt que les opérations de $\pi$ \inserted{(Bien entendu, $I'$ opère trivialement \unsure{sur} $I'_{\unsure{ab}}$, donc c'est le groupe $\pi/I'$ qui opère sur $I'_{\unsure{ab}}$ et à fortiori sur $I'_t$.)} Ceci posé, on constate que $I/I'$ opère trivialement sur (3.16.1), et que l'opération de $\pi/I\simeq\pi_0$ sur $I_t$ \unsure{ainsi}}

\deleted{\unsure{Or}} l'isomorphisme \deleted{de $\pi$} (3.16.1) respecte les opérations de $\pi$ sur les deux mbs, où l'opération sur le deuxième mb. \deleted{est} se fait via $\pi_0$ comme expliqué déjà plus haut.

Ceci en fait exprime le fait que l'isom canonique (\unsure{2.15.1|3.15.1}) est compatible aux ``transport de structure''.

En particulier, on trouve que $I\deleted{/I'}$ opère trivialement sur $I'_t$, en d'autres termes

\uline{Corollaire} $3.\unsure{16}$. Pour tt ss-groupe \inserted{ouvert} invariant $I'$ \deleted{d'indice fini} de $I$, l'extension naturelle de $I/I'$ par le quotient $I'_t\simeq\prod_{\ell\ne p}\mathbb Z_\ell(1)$ de $I'$ est une extension \uline{centrale}.
\margin{16}
